\scanpage{17}
\uline{Cohomologie à supports compacts} \deleted{\unsure{de}} de $C,E,\widehat E$ ....

La cohomologie à \unsure{support compact|supports compacts} des $E,\widehat E$ se déduit simplement par dualité des \deleted{celles} \inserted{\unsure{la Ch. usuelle|les Ch. usuelles}} des $E,\widehat E$ lorsque $X$ est lisse ... Celle de $\widehat C$ est identique à sa cohomologie usuelle (car $\widehat C$ est propre...).

Reste à déterminer la ch. à supports compacts de $C$.

\unsure{Car.|Or.}
\[
\deleted{H_c^*(C')\longleftarrow H_c^{\unsure{p|*}}(\unsure{\cdot|\ldots})}
\]

\diagram{\begin{tikzcd}[ampersand replacement=\&,column sep=tiny,row sep=small]
\cdots\arrow[r,dash] \& H^i(\widehat C)\arrow[r,dash]\arrow[d,phantom,"\rotatebox{90}{$\simeq$}" description] \& H^i(X)\arrow[r,dash] \& H_c^{i+1}(C)\arrow[r,dash] \& H^{i+1}(\widehat C)\arrow[r,dash]\arrow[d,phantom,"\rotatebox{90}{$\simeq$}" description] \& H^{i+1}(X)\\
{} \& \substack{H^{i-2}(X)(-1)\\ \text{si }i\ne0} \& {} \& {} \& H^{i-1}(X)(-1) \& {}
\end{tikzcd}}

Comparant avec la suite exacte analogue pour déterminer \inserted{la ch. de} $E$ \inserted{\unsure{ou celle donnant} $H_x^*(\widehat C)$ \unsure{pr|par} $\widehat E\subset\widehat C$,} on est amené à prévoir $H^i(E)\overset{\text{si }i\ne0}{\simeq}H_c^{i+1}(C)\mathrel{\underset{\unsure{a|\sim}}{\simeq}}H_{\deleted{\unsure{c|x}}x}^{i+1}(C)$, \deleted{\unsure{et}} \deleted{\unsure{on a en effet}, \unsure{grâce à l'inclusion} $E\subset C$,}
\[
\boxed{H_x^*(C)\simeq H_c^*(C)\quad ??}
\]
\[
\deleted{H_x^i(C)=H_c^i(C)=}
\]
% archive sheet number: 16
